\subsection{Differential forms, exterior differentiation}

Recall that, in this issue and the following ones, we assume either that K has characteristic 0, or that the varieties under consideration are locally of finite dimension.

8.3.1. Let X be a manifold of class \(\mathbf{C}^r\) and let F be a vector bundle of class \(\mathbf{C}^k\) with base X, with \(k + 1 \in \mathbb{N}_{\mathbb{K}}\) and \(k + 1 \leqslant r\) . Let \(p\) be an integer \(\geqslant 0\) and let \(\operatorname{Alt}^p(\mathrm{T}(\mathrm{X}); \mathrm{F})\) be the vector bundle of class \(\mathbf{C}^k\) of \(p\)-alternating multilinear maps from \(\mathrm{T}(\mathrm{X})\) into F (7.8.1 and Errata to no. 7.8). A section of this bundle over an open set U of X is called an alternating differential form (or simply differential form) of degree \(p\) on U, with values in F. Those of these sections that are of class \(\mathbf{C}^k\) form a \(\mathcal{C}^k(\mathrm{U})\)-module, denoted \(^k\Omega^p(\mathrm{U}; \mathrm{F})\). If E is a Banach space, one writes \(^k\Omega^p(\mathrm{U}; \mathrm{E})\) instead of \(^k\Omega^p(\mathrm{U}; \mathrm{E}_{\mathrm{x}})\) and one speaks of differential forms with values in E. In these notations, one sometimes omits \(k\) (resp. E) when it is equal to \(r - 1\) (resp. K).

A differential form of degree 0 (resp. 1) with scalar values is a function (resp. a field of covectors).

Let \(\varphi\colon\mathbf{Y}\to\mathbf{X}\) be a morphism of manifolds of class \(\mathbf{C}^{r}\) , and let \(\omega\in{}^{k}\Omega^{p}(\mathbf{X};\mathbf{F})\). There then exists a form \(\varphi^{*}\omega\in{}^{k}\Omega^{p}(\mathbf{Y};\varphi^{*}\mathbf{F})\) and only one such that

\[
(\varphi^ {*} \omega) _ {y} (v _ {1}, \dots , v _ {p}) = \omega_ {\varphi (y)} (\mathrm{T} _ {y} (\varphi). v _ {1}, \dots , \mathrm{T} _ {y} (\varphi). v _ {p})\tag{1}
\]

for all \(y \in Y\) and every family \(v_{1}, \ldots, v_{p}\) of elements of \(\mathrm{T}_{y}(\mathrm{Y})\) .

If \(\psi: Z \to Y\) is a morphism of manifolds of class \(C^r\) , one has

\[
(\varphi \circ \psi) ^ {*} \omega = \psi^ {*} (\varphi^ {*} \omega).\tag{2}
\]

When Y is a submanifold of X and \(\varphi\) is the canonical injection, one sometimes writes \(\omega|Y\) instead of \(\varphi^{*}\omega\) and one says that \(\omega|Y\) is induced on Y by \(\omega\) .

8.3.2. The definitions and results of no. 7.8 apply to differential forms. In particular, a pairing of vector bundles \(F' \times_{\mathbf{X}} F'' \to F\) gives rise to an exterior product \(^k\Omega^{p'}(\mathrm{U};\mathrm{F}') \times ^k\Omega^{p''}(\mathrm{U};\mathrm{F}'') \to ^k\Omega^p(\mathrm{U};\mathrm{F})\) , where \(p = p' + p''\) (7.8.2); it is denoted \((\omega', \omega'') \mapsto \omega' \wedge \omega''\) . The direct sum of the \(^k\Omega^p(\mathrm{U};\mathrm{K})\) for \(p \geqslant 0\) is thus endowed with a structure of associative and alternating graded algebra (A, III, p. 53).

Let \(\xi\) be a vector field on X and \(\omega\) a differential form of degree \(p \geqslant 1\) on X, with values in a vector bundle F; the interior product of \(\xi\) and \(\omega\) is the differential form \(i(\xi, \omega)\) of degree \(p - 1\) on X, with values in F, defined by

\[
i (\xi , \omega) _ {x} (v _ {1}, \dots , v _ {p - 1}) = \omega_ {x} (\xi (x), v _ {1}, \dots , v _ {p - 1})\tag{3}
\]

for all \(x \in X\) and every family \(v_{1}, \ldots, v_{p-1}\) of elements of \(\mathrm{T}_{x}(\mathrm{X})\) . If \(\omega\) is a differential form of degree 0, we set \(i(\xi, \omega) = 0\) . We also write \(i(\xi)\omega\) or \(i_{\xi}\omega\) in place of \(i(\xi, \omega)\) . When \(\mathrm{F} = \mathrm{K}_{\mathrm{X}}\) , the definition given above coincides with that of 7.8.4.

Let \(\mathbf{F}' \times_{\mathbf{x}} \mathbf{F}'' \to \mathbf{F}\) a pairing of vector bundles. For \(\omega' \in {}^k \Omega^{p'}(\mathbf{U}; \mathbf{F}')\) and \(\omega'' \in {}^k \Omega^{p''}(\mathbf{U}; \mathbf{F}'')\) , one has

\[
i (\xi) \left(\omega^ {\prime} \wedge \omega^ {\prime \prime}\right) = \left(i (\xi) \omega^ {\prime}\right) \wedge \omega^ {\prime \prime} + (- 1) ^ {p ^ {\prime}} \omega^ {\prime} \wedge i (\xi) \omega^ {\prime \prime}.\tag{4}
\]

8.3.3. Let \((u^{1},\ldots,u^{n})\) be a coordinate system in the open set U of X. Every element \(\omega\) of \({}^{k}\Omega^{p}(\mathrm{U};\mathrm{F})\) is written uniquely

\[
\omega = \sum_ {i _ {1} <   \dots <   i _ {p}} f _ {i _ {1}, \dots , i _ {p}}. d u ^ {i _ {1}} \wedge \dots \wedge d u ^ {i _ {p}}
\]

where the \(f_{i_1,\ldots,i_p}\) are sections of class \(\mathbf{C}^k\) of \(\mathbf{F}\) on \(\mathbf{U}\) ; in this formula, the sign \(\wedge\) is relative to the canonical pairing \(\mathbf{K}_{\mathbf{x}} \times_{\mathbf{x}} \mathbf{K}_{\mathbf{x}} \to \mathbf{K}_{\mathbf{x}}\) and the point denoting the product is relative to the canonical pairing \(\mathbf{F} \times_{\mathbf{x}} \mathbf{K}_{\mathbf{x}} \to \mathbf{F}\) .

8.3.4. Let \(p\) be an integer \(\geqslant 0\) and \(r\) an element of \(\mathbf{N}_{\kappa}\) . Let \(E\) and \(F\) be Banach spaces, \(U\) an open subset of \(E\) and \(\alpha \in {}^r\Omega^p(U;F)\) . Let \(\tilde{\alpha} \in \mathcal{C}^r (U; \operatorname{Alt}^p(E;F))\) be the element corresponding to \(\alpha\) . If \(x \in U\) , the differential \(d_x\tilde{\alpha}\) of \(\tilde{\alpha}\) at \(x\) (5.5.6) is an element of \(\mathcal{L}(E; \operatorname{Alt}^p(E;F))\) ; if \(t \in E\) , one has \((d_x\tilde{\alpha})(t) \in \operatorname{Alt}^p(E;F)\) , and if \(t_1, \ldots, t_p\) are in \(E\) , one has \((d_x\tilde{\alpha})(t)(t_1, \ldots, t_p) \in F\) . There exists one and only one element \(d\alpha\) of \({}^{r-1}\Omega^{p+1}(U;F)\) such that one has

\[
(d \alpha) _ {x} (t _ {0}, \dots , t _ {p}) = \sum_ {i = 0} ^ {p} (- 1) ^ {i} (d _ {x} \tilde {\alpha}) (t _ {i}) (t _ {0}, \dots , \hat {t} _ {i}, \dots , t _ {p}) ^ {(1)}\tag{5}
\]

for all \(x \in U\) and \(t_{0}, \ldots, t_{p}\) in E.

8.3.5. Let \(p\) be an integer \(\geqslant 0\) and \(r\) an element of \(\mathbf{N}_{\mathbf{K}}\) . Let \(X\) be a manifold of class \(\mathbf{C}^{r+1}\) , \(F\) a Banach space, and \(\omega \in {}^r\Omega^p(X;F)\). There exists a differential form \(\pi \in {}^{r-1}\Omega^{p+1}(X;F)\) and only one such that, for every chart \(c = (U,\varphi,E)\) of \(X\) , if \(\omega_c\) is the differential form on \(\varphi(U)\) such that \(\omega|U = \varphi^*(\omega_c)\) , one has \(\pi_c = d\omega_c\) , where \(d\omega_c\) is defined by formula (5) of 8.3.4.

The differential form \(\pi\) is called the exterior differential of \(\omega\) ; it is denoted \(d\omega\) . It has the following properties:

\begin{enumerate}
\def\labelenumi{(\arabic{enumi})}
\setcounter{enumi}{5}
\item
  If \(\psi: Y \to X\) is a morphism of varieties of class \(C^{r+1}\) and if \(\omega \in {}^r\Omega^n(X; F)\) , one has \(\psi^*(d\omega) = d(\psi^*\omega)\) ; in particular, \(d\) commutes with the operation of restriction to a submanifold.
\item
  For \(p = 0\) , the map \(d: ^r\Omega^0(\mathbf{X}; \mathbf{F}) \to {}^{r-1}\Omega^1(\mathbf{X}; \mathbf{F})\) coincides with the one defined in 8.2.2.
\item
  If \(r \geqslant 2\) and if \(\omega \in {}^r\Omega^p(\mathbf{X};\mathbf{F})\) , one has \(d(d\omega) = 0\) . \(^{(2)}\)
\item
  For every continuous linear map \(u: \mathbf{F} \to \mathbf{F}'\) of Banach spaces, we have \(d(u(\omega)) = u(d(\omega))\) for \(\omega \in {}^{r}\Omega^{p}(\mathbf{U}; \mathbf{F})\) .
\item
  For every continuous bilinear mapping \(F' \times F'' \rightarrow F\) of Banach spaces, we have
\end{enumerate}

\[
d \left(\omega^ {\prime} \wedge \omega^ {\prime \prime}\right) = \left(d \omega^ {\prime}\right) \wedge \omega^ {\prime \prime} + (- 1) ^ {p ^ {\prime}} \omega^ {\prime} \wedge d \omega^ {\prime \prime}
\]

for \(\omega' \in {}^r\Omega^{p'}(\mathbf{X}; \mathbf{F}')\) and \(\omega'' \in {}^r\Omega^{p''}(\mathbf{X}; \mathbf{F}'')\) .


8.3.6. Let \((u^{1},\ldots ,u^{n})\) be a coordinate system of X in U. Let

\[
\omega = \sum_ {i _ {1} <   \dots <   i _ {p}} f _ {i _ {1}, \dots , i _ {p}}. d u ^ {i _ {1}} \wedge \dots \wedge d u ^ {i _ {p}}
\]

\(^{1}\) The sign \^{} indicates that the symbol it surmounts is to be omitted (cf.~7.8.4).

\(^{2}\) If \(\omega\in^{1}\Omega^{p}(X;F)\) and if \(d\omega\in^{1}\Omega^{p+1}(X;F)\) , we still have \(d(d\omega)=0\) .

a differential form on U (with \(f_{i_1,\ldots,i_p} \in \mathcal{C}^r(\mathrm{U};\mathrm{F})\) ). We then have

\[
d \omega = \sum_ {i _ {1} <   \dots <   i _ {p}} d f _ {i _ {1}, \dots , i _ {p}} \wedge d u ^ {i _ {1}} \wedge \dots \wedge d u ^ {i _ {p}}.\tag{11}
\]

8.3.7. Suppose K has characteristic zero. Let \(\omega \in {}^r\Omega^p (\mathbf{X};\mathbf{F})\) be a differential form of degree \(p\geqslant 1\) such that \(d\omega = 0\) . For every \(x\in \mathbf{X}\) there exists an open neighborhood U of \(x\) and a differential form \(\pi \in {}^r\Omega^{p - 1}(\mathrm{U};\mathrm{F})\) such that \(d\pi = \omega\) on U.

